\section{Forecast-optimal smoothness by chronological cross-validation}
\label{sec:forecast_cv}

\subsection{From a fitted trend to a prospective forecast}

The parameter of interest is the normalized smoothness $S$, not the
unbounded penalty $\lambda$. Define the continuous family of
smoothing matrices, including their limiting endpoint, by
\begin{equation}
\label{eq:sindexed_smoothing_matrix}
H(S)=
\begin{cases}
\bigl(I+\lambda(S)D_d^\top D_d\bigr)^{-1},&0\le S<1,\\
P_{\ker(D_d)},&S=1,
\end{cases}
\end{equation}
where $P_{\ker(D_d)}$ is the orthogonal projection onto the
null space of the finite-difference operator. Computing
$H(S)$ may require inverting the monotone map $S(\lambda)$,
but the domain on which cross-validation operates is the
\emph{full bounded interval} $[0,1]$.

A fitted trend is not automatically a forecast. We specify
a continuation rule that imposes zero future $d$th
differences, producing the fixed linear operator
\begin{equation}
\label{eq:continuation_operator_space}
G_{d,h}\in\mathbb{R}^{h\times L}.
\end{equation}
The corresponding prospective forecast from origin $t$ is
\begin{equation}
\label{eq:forecast_operator}
\widehat z_t(S)=G_{d,h}H(S)x_t.
\end{equation}
In terms of the final backward differences of the
historically fitted trend,
\begin{equation}
\label{eq:continuation}
\widehat\tau_{t+k\mid t}(S)
=\sum_{j=0}^{d-1}\binom{k+j-1}{j}
\nabla^j\bigl[H(S)x_t\bigr]_{L},
\qquad k=1,\ldots,h.
\end{equation}
For $d=1$, $d=2$, and $d=3$, these continuations are constant,
linear, and quadratic in the forecast step $k$, respectively.
More generally, fixing $d$ fixes the maximum degree $d-1$,
while changing $S$ changes the smoothed endpoint values and
differences, hence the coefficients of the forecast polynomial.
The method selects among \emph{future extrapolations}, not
merely among different post-hoc representations of the past.
Selecting $d$ itself is a separate model-selection decision;
it is held fixed within the baseline CV comparison.

\subsection{The forecast-MSE surface in the smoothness coordinate}

For an \emph{historical} forecast origin $t$, let
\begin{equation}
\label{eq:future_observation_block}
z_t=(y_{t+1},\ldots,y_{t+h})^\top
\end{equation}
be the subsequent realized validation block. The loss assigned
to a candidate smoothness $S\in[0,1]$ is
\begin{equation}
\label{eq:origin_loss_smoothness}
F_t(S)
=\frac1h\left\|z_t-G_{d,h}H(S)x_t\right\|_2^2.
\end{equation}
Expanding this expression makes the dependence on the
smoothness-indexed smoothing matrix explicit:
\begin{equation}
\label{eq:origin_loss_quadratic}
\begin{aligned}
h F_t(S)
={}&z_t^\top z_t
-2z_t^\top G_{d,h}H(S)x_t\\
&+x_t^\top H(S)G_{d,h}^\top
G_{d,h}H(S)x_t.
\end{aligned}
\end{equation}
The matrix $H(S)$ is symmetric. Unlike ordinary fitting
criteria, \cref{eq:origin_loss_smoothness} scores observations
\emph{after} the training endpoint; the in-sample residual
$\|x_t-H(S)x_t\|^2$ does not appear in the criterion.

For derivative calculations it is sometimes convenient to
use the algebraically equivalent auxiliary function
\begin{equation}
\label{eq:origin_loss}
f_t(\lambda)=F_t(S(\lambda)).
\end{equation}
This equality is a coordinate change, not a second
optimization problem or a requirement to search a truncated
range of $\lambda$.

\subsection{Chronological, horizon-matched forecast cross-validation}

At final forecast origin $T$, select historical origins
$t_1,\ldots,t_M$ whose training windows are available and
whose validation blocks are already observed:
\begin{equation}
\label{eq:chronological_origin_restriction}
t_m+h\le T,\qquad m=1,\ldots,M.
\end{equation}
For fixed $(d,L,h)$, the proposed CV surface is
\begin{equation}
\label{eq:F_S}
F^{\mathrm{pool}}_{T,d,L,h}(S)
=\frac1M\sum_{m=1}^M F_{t_m}(S),
\qquad S\in[0,1].
\end{equation}
It produces the forecast-optimal smoothness estimate
\begin{equation}
\label{eq:Sstar}
\widehat S^{\mathrm{FCV}}_{T,d,L,h}
\in\operatorname*{arg\,min}_{S\in[0,1]}
F^{\mathrm{pool}}_{T,d,L,h}(S).
\end{equation}
The selection space includes both exact limiting models,
$S=0$ and $S=1$, not just interior grid values.
When the objective has several minima, the global
minimizer is determined by the lowest historical predictive
MSE including the endpoints; a deterministic smallest-$S$
tie convention is used if several minima have the same
recorded objective value.

Both the forecasting horizon and the continuation operator
are essential parts of the criterion. Minimizing one-step
forecast error and deploying the resulting smoothness for
a longer horizon is a different procedure. The index makes
this comparison interpretable on $[0,1]$, while its
one-to-one transformation from $\lambda$ preserves the
fitted-trend family and any exact minimizing set.

\subsection{Final refit and untouched future forecast}

At outer origin $T$, all validation-stage trends are
discarded. Using the smoothness chosen from historical
blocks, the current trend is refitted on the newest
$L$ observations:
\begin{equation}
\label{eq:final_refit}
\widehat\tau_T
=H\bigl(\widehat S^{\mathrm{FCV}}_{T,d,L,h}\bigr)
y_{T-L+1:T}.
\end{equation}
The operational forecast is
\begin{equation}
\label{eq:final_operational_forecast}
\widehat y_{T+1:T+h\mid T}
=G_{d,h}\widehat\tau_T.
\end{equation}
No observation after $T$ enters either the CV decision
or the final fit. Historical forecast losses, not the
historical fitted trend paths, are averaged.

\subsection{Relationship to established cross-validation}

Ordinary leave-one-out CV and GCV primarily measure
interpolation or fitted-signal error for a smoother.
The present criterion uses chronological, genuinely
future-block MSE from a declared polynomial continuation.
This does not invent the general principle of predictive
cross-validation: Hart's time-series CV
\cite{hart1994tscv} is an important precedent.
The specific object investigated here is direct
$h$-step forecast-MSE selection for the
finite-difference PLS trend over normalized $S\in[0,1]$.
Time-adaptive decisions using historical minima, considered
later, are optional and not part of this definition.
